% part: intuitionistic-logic
% chapter: propositions-as-types
% section: proofs-to-terms

\documentclass[../../../include/open-logic-section]{subfiles}

\begin{document}

\olfileid{int}{pty}{pt}

\olsection{\usetoken{P}{derivation} په ثبوت ترمونو بدلول}

د فطري استنتاج !!{proof}s په ثبوت ترمونو بدلول، يعنې
د \Log{N2i}~!!{proof}s د \Log{tN2i}~په !!{proof}s
اړول، ډېر نېغ بهير دے۔ يوازې دومره کوو چې په
!!{derivation} کښې د هر سېکوېنټ په ښي‌لاس اړخ کښې
د هر !!{formula} کيڼې خوا ته يو ثبوت ترم وليکو؛ په
پايله کښې د $\Gamma \Sequent M : !A$ په بڼه
سېکوېنټونه ترلاسه کېږي۔ بيا وايو چې $M$ په
سياق~$\Gamma$ کښې د~$!A$ \emph{شاهد} دے۔ کوم ثبوت
ترم ليکل پکار دي، دا د \Log{tN2i} قواعد په بشپړ
ډول ټاکي۔

\begin{prop}
په \Log{N2i} کښې د $\Gamma \Sequent !A$ هر
!!{proof}~$\delta$ ته په \Log{tN2i} کښې د
$\Gamma \Sequent N : !A$ !!a{proof}~$\delta'$
ورسره برابر وي۔ ثبوت ترم $N$ په يکتا ډول د~$\delta$
له مخې ټاکل کېږي۔
\end{prop}

\begin{proof}
د~$\delta$ په لوړوالي استقرا کوو۔ که $\delta$ د
$x:!A \Sequent !A$ اکسيوم وي، نو $\delta'$ د
$x:!A \Sequent x:!A$ اکسيوم دے۔ که $\delta$ په
يوه استنتاجي قاعده پاې ته رسېږي، استقرايي فرضيه
هرې مقدمې ته ثبوت ترم ورکوي او د دغو مقدمو سره
برابر د سېکوېنټونو \Log{tN2i}-!!{proof}s ورکوي چې
دغه ثبوت ترمونه پکې شامل وي۔ د \Log{tN2i} اړونده
قاعده کاروو؛ اړوند ثبوت ترم د \Log{tN2i} استنتاجي
قاعده ټاکي۔
\end{proof}

\begin{ex}
  په \Log{N1i} کښې د $(!A \land !B) \lif (!B \land !A)$
  دا ډېر ساده !!{proof} وګورئ:
\[
  \AxiomC{$\Discharge{!A\land !B}{x}$}
  \RightLabel{\Elim\land[2]}
  \UnaryInfC{$!B$}
  \AxiomC{$\Discharge{!A\land !B}{x}$}
  \RightLabel{\Elim\land[1]}
  \UnaryInfC{$!A$}
  \RightLabel{\Intro\land}
  \BinaryInfC{$!B \land !A$}
  \DischargeRule{\Intro\lif}{x}
  \UnaryInfC{$(!A \land !B) \lif (!B \land !A)$}
  \DisplayProof
  \]
  په \Log{N2i} کښې ئې بڼه داسې ده:
\[
  \Axiom$x : !A\land !B \fCenter !A\land !B$
  \RightLabel{\Elim\land[2]}
  \UnaryInf$x : !A\land !B \fCenter !B$
  \Axiom$x : !A\land !B \fCenter !A\land !B$
  \RightLabel{\Elim\land[1]}
  \UnaryInf$x : !A\land !B \fCenter !A$
  \RightLabel{\Intro\land}
  \BinaryInf$x : !A\land !B \fCenter !B \land !A$
  \RightLabel{\Intro\lif}
  \UnaryInf$\fCenter (!A \land !B) \lif (!B \land !A)$
  \DisplayProof
  \]
  ثبوت ترمونه له پاسه ښکته ټاکو۔ د اکسيومونو ثبوت
  ترمونه هماغه د سياق متغير~$x$ دے۔ بيا له
  $\Elim\land[i]$ سره تړلي ثبوت ترمونه
  $\proj{2}{x}$ او~$\proj{1}{x}$ ټاکل کېږي۔ کله چې
  \Intro{\land} کاروو، دغه دوه ثبوت ترمونه يوځاے
  کوو: $\pair{\proj{2}{x}}{\proj{1}{x}}$۔ په پاې
  کښې د \Intro{\lif} قاعده يو $\lambd$-تجريد
  کاروي،~$x$ تړي او ثبتوي چې $x$ د~$!A \land !B$
  فرض نښه وه: $\lambd[\typeof{x}{!A \land
  !B}][\pair{\proj{2}{x}}{\proj{1}{x}}]$۔ نو په
  \Log{tN2i} کښې !!{proof} داسې کېږي:
\[
  \Axiom$x : !A\land !B \fCenter x: !A\land !B$
  \RightLabel{\Elim\land[2]}
  \UnaryInf$x : !A\land !B \fCenter \proj{2}{x} : !B$
  \Axiom$x : !A\land !B \fCenter x: !A\land !B$
  \RightLabel{\Elim\land[1]}
  \UnaryInf$x : !A\land !B \fCenter \proj{1}{x} : !A$
  \RightLabel{\Intro\land}
  \BinaryInf$x : !A\land !B \fCenter \pair{\proj{2}{x}}{\proj{1}{x}} : 
    !B \land !A$
  \RightLabel{\Intro\lif}
  \UnaryInf$\fCenter \lambd[\typeof{x}{!A \land
  !B}][\pair{\proj{2}{x}}{\proj{1}{x}}] : (!A \land !B) \lif (!B \land !A)$
  \DisplayProof
  \]
\end{ex}

\end{document}
